% D. Metodología desarrollada de forma sucesiva (Guía apdo. 3.D)
% Extensión orientativa: 4-5 páginas. Narrar las fases EN ORDEN, incluidos los problemas
% encontrados y cómo se resolvieron: es lo que evidencia el trabajo realizado.
% Fuente de verdad del avance: docs/ROADMAP.md
\chapter{Metodología}
\label{cap:metodologia}

El desarrollo se organizó en fases sucesivas, cada una condicionada por el resultado de la
anterior. Este capítulo describe esa progresión, incluidas las incidencias encontradas y
las soluciones adoptadas.

\section{Fase 1: definición del escenario y de la arquitectura}

Se seleccionó como activo de prueba el compresor de un sistema de refrigeración doméstico y
se definieron las tres modalidades de sensado a integrar. Sobre esa base se estableció una
arquitectura de dos niveles: un nodo de captura en el borde y un concentrador local que
recibe la telemetría y construye el conjunto de datos.

\section{Fase 2: integración del nodo de sensado}

La integración se abordó de forma incremental, incorporando y validando un sensor cada vez
antes de combinarlos en un único ciclo de adquisición. El orden seguido fue: acelerómetro y
giroscopio sobre bus \textit{Inter-Integrated Circuit} (I\textsuperscript{2}C), sonda de
temperatura sobre bus 1-Wire y micrófono digital sobre bus \textit{Inter-IC Sound} (I2S).

\section{Fase 3: estabilización frente a la vibración}

Durante las primeras pruebas con el compresor en funcionamiento se detectó un fallo
recurrente: el bus I\textsuperscript{2}C quedaba bloqueado y el bucle principal del
microcontrolador se detenía indefinidamente. El origen se localizó en microcortes del
cableado del acelerómetro provocados por la propia vibración del activo, precisamente la
magnitud que el sistema debe medir.

La solución adoptada fue una rutina de detección y recuperación por software: en cada ciclo
se comprueba el reconocimiento del dispositivo en su dirección de bus y, si no responde, se
reinicia el bus y se reinicializa el sensor sin reiniciar el microcontrolador. Se descartó
de forma deliberada la alternativa de reiniciar el nodo, dado que un reinicio interrumpe la
continuidad temporal de la señal, que es la información que el análisis posterior explota.

\section{Fase 4: conectividad y almacenamiento}

Al preparar la primera campaña de captura prolongada se identificó una limitación en el
esquema de adquisición que habría invalidado el análisis previsto. El bucle de medida
operaba a una cadencia de \mbox{1 Hz}, suficiente para las magnitudes térmicas pero dos
órdenes de magnitud por debajo de lo que exige la vibración del activo, cuya componente
fundamental se sitúa en torno a los \mbox{48 Hz}. La consecuencia no era una pérdida de
precisión, sino la imposibilidad de todo análisis frecuencial: la señal se replegaba sobre la
banda observable.

La corrección consistió en desacoplar las dos escalas temporales del problema mediante dos
canales de adquisición con cadencias independientes, según se detalla en el
capítulo~\ref{cap:diseno}. El formato completo de las tramas de ambos canales, los valores
centinela que emplean y los procedimientos de aprovisionamiento y verificación del
concentrador se recogen en el anexo~\ref{anexo:software}. El bucle principal se reformuló además para no depender de
esperas activas prolongadas, y la conversión de la sonda de temperatura pasó a solicitarse de
forma asíncrona, dado que sus aproximadamente \mbox{750 ms} de conversión detenían el resto
de la adquisición.

Procede señalar que la limitación se detectó antes de lanzar la campaña de captura y no
después. De haberse iniciado la adquisición con el esquema anterior, las horas de medida
registradas no habrían contenido información espectral aprovechable y la campaña habría
debido repetirse por completo.

La validación del esquema de ráfagas se realizó en dos etapas deliberadamente separadas.
Las rutinas de extracción de características se implementaron sin dependencias del entorno
de ejecución del microcontrolador, lo que permitió verificarlas en un computador frente a
señales sintéticas de propiedades conocidas analíticamente. Solo después se comprobó su
comportamiento en la placa. El motivo de este orden es práctico: depurar operaciones
matemáticas a través del puerto serie del microcontrolador resulta desproporcionadamente
lento frente a hacerlo en un entorno de compilación local.

La verificación sobre la placa confirmó los dos parámetros críticos del esquema. La duración
real de la captura se mantuvo en \mbox{1024 ms} sin desviación apreciable en todas las
ráfagas observadas, lo que acredita que la cadencia de muestreo se sostiene y que el eje de
frecuencias resultante es válido. El cálculo completo de las características de los tres
ejes y de la ventana acústica consumió entre \mbox{29 ms} y \mbox{30 ms}, margen holgado
frente al periodo de \mbox{30 s} entre ráfagas. Se comprobó igualmente que, con el activo en
marcha, la frecuencia dominante coincidía en los tres ejes en torno a \mbox{49 Hz}, valor
compatible con el régimen de un motor de inducción de dos polos alimentado a \mbox{50 Hz} una
vez descontado el deslizamiento. Ese pico desaparecía con el compresor detenido, lo que
permitió atribuirlo al giro del motor y no a un acoplamiento eléctrico de la red.

\section{Fase 5: despliegue del concentrador}

El concentrador se implantó sobre una Raspberry~Pi que aloja el intermediario de mensajería
y el registrador. Durante la puesta a punto se constató que el equipo que alojaba el
intermediario cambió de subred en cuatro ocasiones, circunstancia problemática porque la
dirección del intermediario es una constante de compilación del firmware y cada cambio
obligaba a reprogramar el nodo.

La solución adoptada consistió en configurar la interfaz inalámbrica del concentrador como
punto de acceso, de forma que el nodo se conecte directamente a él sin encaminador
intermedio. El gestor de red asigna en ese modo una dirección fija, con lo que la constante
del firmware deja de depender de infraestructura ajena. La decisión aporta además dos
ventajas: el conjunto resulta portátil, al no requerir más que alimentación eléctrica, y la
red queda aislada, lo que elimina el riesgo de que un tercero publique tramas en los temas de
telemetría y contamine el conjunto de datos. Se asume como contrapartida que el concentrador
pierde su acceso a Internet por la interfaz inalámbrica, dado que esta no puede operar
simultáneamente como cliente y como punto de acceso; la gestión se mantiene por cable.

El aprovisionamiento se automatizó separando las operaciones en fases, por una razón de
seguridad operativa: la activación del punto de acceso interrumpe la red por la que llega la
sesión de administración remota. Todas las comprobaciones se realizan antes de modificar la
configuración del sistema, la activación del punto de acceso constituye la última operación y
se lanza desacoplada de la sesión para que se complete aunque la conexión se pierda.

El equipo empleado como concentrador contaba ya con un intermediario configurado para otros
usos. Se adoptó como criterio la intervención mínima: el procedimiento examina la
configuración existente y, cuando esta ya satisface los requisitos, no la modifica; cuando la
configuración ajena restringe el servicio a la interfaz local, se limita a informar sin
alterar ficheros de terceros. Ante cualquier fallo, restaura el estado anterior antes de
abortar. El fundamento es que dejar un servicio inoperativo en un equipo compartido resulta
más perjudicial que no haber intervenido.

La cadena completa se verificó eslabón por eslabón siguiendo el recorrido del dato: punto de
acceso activo, nodo asociado a la red, intermediario aceptando conexiones, tramas de ambos
canales recibidas, registrador en ejecución y ficheros de datos incrementando su número de
filas. Esta última comprobación se realiza contrastando el recuento de filas en dos instantes
separados, y no verificando la mera existencia del fichero, porque un fichero presente que no
crece constituye el modo de fallo más fácil de pasar por alto. Se comprobó asimismo que el
conjunto se restablece por completo tras un corte de alimentación sin intervención manual.

Las fases experimentales y los resultados obtenidos a partir de las campañas de captura
se recogen y analizan en detalle en el capítulo~\ref{cap:resultados}.

El preprocesado comprende dos filtrados de naturaleza distinta. El primero descarta los
valores centinela descritos en el capítulo~\ref{cap:diseno}, cuya inclusión distorsionaría
cualquier estadístico: un único valor de \mbox{$-127$ \textdegree C} arrastra la media de
temperatura de una jornada completa.

El segundo es consecuencia directa de las limitaciones del montaje documentadas en el
anexo~\ref{anexo:hardware}. Dado que se asume una tasa de fallo de lectura apreciable en
lugar de intervenir sobre el cableado, el conjunto de datos contendrá ráfagas de calidad
desigual, y el protocolo debe discriminarlas en lugar de tratarlas por igual. Los contadores
de salud que el nodo publica hacen posible ese filtrado sobre criterios objetivos: la
duración real de la captura permite descartar las ráfagas cuya cadencia no se sostuvo, y el
número de reintentos de cada ráfaga califica su calidad individual. Se establece por tanto
como criterio metodológico que el filtrado del conjunto de datos se documente de forma
explícita, indicando qué fracción de las ráfagas se descartó y con qué umbral, dado que esa
fracción forma parte del resultado y no de su preparación.

Sobre las ráfagas retenidas se calculan las características derivadas que requieren
historia y que el nodo no puede obtener de una ventana aislada: el gradiente térmico, el
diferencial entre la temperatura de la superficie del compresor y la del ambiente, la
evolución de los armónicos respecto de la componente fundamental y la duración de los ciclos
de marcha y parada, deducida del canal lento.

El diferencial térmico merece una mención por su función en el diseño. Emplear la diferencia
entre la temperatura del activo y la del ambiente, en lugar de la temperatura absoluta,
elimina de la formulación la variación estacional: un compresor en estado nominal mantiene un
salto térmico semejante sobre el ambiente con independencia de la temperatura exterior. La
alternativa, un modelo que persiguiese esa deriva ajustándose de forma continua, resulta
inviable porque no puede distinguir una deriva ambiental lenta de una degradación igualmente
lenta. Se prefiere por ello eliminar la variable de confusión a compensarla.

La validación exige disponer de referencia de verdad, que en este contexto solo puede
obtenerse provocando los fallos de forma deliberada. Se establecieron cuatro modos, elegidos
por ser reversibles y por afectar a modalidades de medida distintas: desequilibrio de masa,
aflojamiento del anclaje, obstrucción de la ventilación y sobrecarga térmica. Los dos
primeros se manifiestan en la vibración, el tercero en las magnitudes térmica y acústica y el
cuarto en la evolución térmica y en el patrón de los ciclos de marcha.

El protocolo de cada campaña impone tres condiciones. La primera es que la campaña de
referencia en estado nominal precede a cualquier fallo inducido, por ser irreversible el
orden inverso. La segunda es que el montaje del sensor no se modifica entre campañas: dado
que la firma de vibración depende del punto y de la rigidez de la fijación, un remontaje
invalida la referencia y obliga a repetirla. La tercera es que toda campaña queda anotada con
su activo, condición, montaje, ficheros e incidencias, sin lo cual un fichero de datos carece
de las condiciones experimentales que le dan sentido.

La comparación se plantea dentro de un mismo activo, con estado nominal y estados de fallo
registrados sobre la misma máquina y el mismo montaje. Esta restricción excluye contrastar el
modelo ajustado sobre un activo frente a las medidas de otro distinto, procedimiento que no
permitiría discernir si una anomalía señalada corresponde al fallo o simplemente a la
diferencia entre máquinas.

El umbral de decisión se calibra sobre la distribución de puntuaciones del conjunto nominal,
lo que fija la tasa de falsos positivos admitida, y la evaluación se expresa mediante la
matriz de confusión frente a cada modo de fallo inducido.

La última fase traslada el sistema a un activo en uso diario, condición que introduce
perturbaciones ausentes en el banco: aperturas de puerta, variación de la carga térmica y
ciclos de desescarche. El objetivo no es mejorar la detección sino comprobar que el
comportamiento normal del activo, en toda su variabilidad, no genera una tasa de aviso
inaceptable.

Las métricas de robustez previstas son el tiempo de servicio sin intervención, la fracción de
ráfagas válidas sobre las esperadas, la evolución de los contadores de salud del nodo y la
tasa de avisos en ausencia de fallo conocido. Se incorpora además la comprobación de que el
veredicto de estado sigue siendo observable en el propio nodo con el concentrador apagado,
dado que la autonomía del extremo constituye la tesis del trabajo y afirmarla sin verificarla
carecería de fundamento.
